\chapter*{Resumen}
\addcontentsline{toc}{chapter}{\footnotesize RESUMEN}

Este Trabajo Fin de Master desarrolla una cadena completa para estimar y explicar la volatilidad implicita de opciones sobre el IBEX a partir de operaciones reales de mercado. La volatilidad implicita no se observa directamente: se obtiene invirtiendo un modelo de valoracion, en este caso Black-76, para encontrar la volatilidad que hace compatible el precio teorico con el precio negociado. Sobre esa base se construye un conjunto de datos depurado, se entrenan varias familias de modelos de regresion y se genera un dashboard orientado a inspeccionar tanto el rendimiento como el comportamiento financiero del modelo.

La metodologia combina un ETL reproducible, validaciones de calidad, feature engineering financiero, particiones temporales sin fuga de informacion, busqueda de hiperparametros mediante k-folds temporales y evaluacion final fuera de muestra. Se comparan regresion lineal, Random Forest, XGBoost, redes neuronales secuenciales y una red neuronal inspirada en tensor-train. Los resultados muestran que los modelos no lineales capturan mejor la superficie de volatilidad que el baseline lineal; en las ejecuciones disponibles, Random Forest alcanza el mejor compromiso cuantitativo, con RMSE de test cercano a 0,0884 y $R^2$ de test cercano a 0,791.

La contribucion principal del trabajo no es solo predictiva, sino explicable. El sistema transforma el modelo final en artefactos de interpretabilidad global y local: SHAP, curvas ICE y ALE, superficies de volatilidad, arboles surrogate, regresion simbolica, vecinos cercanos y diagnosticos de residuos. Estos elementos permiten analizar que variables sostienen una prediccion, donde falla el modelo y hasta que punto las aproximaciones interpretables son fieles al estimador principal.
